\documentclass[10pt, aspectratio=169]{beamer}
% Preámbulo modular para presentaciones Beamer (Modelación Estadística)
% Diseño y paleta institucional del Tecnológico de Monterrey

\documentclass[aspectratio=169,xcolor={dvipsnames,table}]{beamer}

\usepackage[utf8]{inputenc}
\usepackage[spanish,mexico]{babel}
\usepackage{amsmath,amssymb,amsfonts,mathtools}
\usepackage{graphicx}
\usepackage{listings}
\usepackage{booktabs}
\usepackage{tikz}
\usetikzlibrary{matrix,arrows,decorations.pathmorphing,shapes.geometric,positioning}

% Paleta Institucional Tecnológico de Monterrey
\definecolor{TecAzul}{HTML}{0039A6}       % Azul TEC: color primario institucional
\definecolor{TecAzulOscuro}{HTML}{1A2E51} % Azul marino oscuro
\definecolor{TecRojo}{HTML}{EC2661}       % Rojo de acento (paleta miTec)
\definecolor{TecGris}{HTML}{646464}       % Gris neutro para comentarios y subtítulos
\definecolor{TecFondoBloque}{HTML}{F4F6F9}% Fondo suave para bloques

% Configuración de Tema Beamer
\usetheme{Boadilla}
\usecolortheme[named=TecAzulOscuro]{structure}
\setbeamercolor{alerted text}{fg=TecRojo}
\setbeamercolor{palette primary}{bg=TecAzulOscuro,fg=white}
\setbeamercolor{palette secondary}{bg=TecAzul,fg=white}
\setbeamercolor{palette tertiary}{bg=TecAzulOscuro!80!black,fg=white}
\setbeamercolor{title}{fg=TecAzulOscuro,bg=TecFondoBloque}
\setbeamercolor{frametitle}{fg=TecAzulOscuro,bg=TecFondoBloque}

% Bloques personalizados elegantes
\setbeamercolor{block title}{bg=TecAzulOscuro,fg=white}
\setbeamercolor{block body}{bg=TecFondoBloque,fg=black}
\setbeamercolor{block title alerted}{bg=TecRojo,fg=white}
\setbeamercolor{block body alerted}{bg=TecRojo!8,fg=black}
\setbeamercolor{block title example}{bg=TecAzul,fg=white}
\setbeamercolor{block body example}{bg=TecAzul!8,fg=black}

% Redondear bordes y añadir sombra sutil
\setbeamertemplate{blocks}[rounded][shadow=true]
\setbeamertemplate{navigation symbols}{}

% Pie de página limpio con número de diapositiva
\setbeamertemplate{footline}{
  \leavevmode%
  \hbox{%
  \begin{beamercolorbox}[wd=.7\paperwidth,ht=2.25ex,dp=1ex,left]{author in head/foot}%
    \hspace*{2ex}\insertshortauthor~~(\insertshortinstitute) \hfill \insertshorttitle
  \end{beamercolorbox}%
  \begin{beamercolorbox}[wd=.3\paperwidth,ht=2.25ex,dp=1ex,right]{date in head/foot}%
    \insertshortdate{}\hspace*{2em}
    \insertframenumber{} / \inserttotalframenumber\hspace*{2ex}
  \end{beamercolorbox}}%
  \vskip0pt%
}

% Configuración de Listings de Python para visualización pedagógica de código
\lstset{
    language=Python,
    basicstyle=\ttfamily\footnotesize,
    keywordstyle=\color{TecAzulOscuro}\bfseries,
    stringstyle=\color{TecRojo},
    commentstyle=\color{TecGris}\itshape,
    showstringspaces=false,
    breaklines=true,
    frame=lines,
    rulecolor=\color{TecAzulOscuro!30},
    backgroundcolor=\color{TecFondoBloque},
    aboveskip=0.5em,
    belowskip=0.5em,
    literate={á}{{\'a}}1 {é}{{\'e}}1 {í}{{\'i}}1 {ó}{{\'o}}1 {ú}{{\'u}}1
             {Á}{{\'A}}1 {É}{{\'E}}1 {Í}{{\'I}}1 {Ó}{{\'O}}1 {Ú}{{\'U}}1
             {ñ}{{\~n}}1 {Ñ}{{\~N}}1 {¿}{{?`}}1 {¡}{{!`}}1
}

% Metadatos institucionales por defecto
\title[Modelación Estadística]{Modelación Estadística para Ciencia de Datos e Ingeniería}
\author[J. Castillo Colmenares]{Juliho David Castillo Colmenares}
\institute[Tec de Monterrey]{Tecnológico de Monterrey \\ Escuela de Ingeniería y Ciencias}
\date{\today}

% Comandos y macros estadísticas compartidas para las presentaciones Beamer (Metropolis)

% Conjuntos numéricos básicos
\newcommand{\R}{\mathbb{R}}
\newcommand{\N}{\mathbb{N}}
\newcommand{\Q}{\mathbb{Q}}
\newcommand{\Z}{\mathbb{Z}}
\newcommand{\set}[1]{\left\{ #1 \right\}}
\newcommand{\abs}[1]{\left|#1\right|}
\newcommand{\norm}[1]{\left\| #1 \right\|}

% Comandos estadísticos y probabilísticos del curso
\newcommand{\Var}{\operatorname{Var}}
\newcommand{\cov}{\operatorname{Cov}}
\newcommand{\corr}{\rho}
\newcommand{\comb}[2]{\begin{pmatrix} #1 \\ #2 \end{pmatrix}}
\newcommand{\s}{\sigma}
\newcommand{\particion}{\mathfrak{P}}
\newcommand{\card}[1]{\#\left( #1 \right)}
\newcommand{\seq}[3]{\set{#1}_{#2}^{#3}}
\newcommand{\E}{\mathbb{E}}
\newcommand{\Prob}{\mathbb{P}}

% Entornos pedagógicos y cajas en español académico natural
\newenvironment{discusion}{%
  \begin{alertblock}{Pregunta de Discusión}%
}{%
  \end{alertblock}%
}

\newenvironment{reflexion}{%
  \begin{alertblock}{Reflexión para el Aula}%
}{%
  \end{alertblock}%
}

\newenvironment{paradoja}{%
  \begin{alertblock}{Paradoja de Exploración}%
}{%
  \end{alertblock}%
}

\newenvironment{motivacion}{%
  \begin{exampleblock}{Motivación: Ciencia de Datos en la Práctica}%
}{%
  \end{exampleblock}%
}

\newenvironment{retoclase}{%
  \begin{block}{Reto Rápido en Equipo}%
}{%
  \end{block}%
}


\title[IC Varianzas y Proporciones]{Sección 06.05: Intervalos de Confianza para Varianzas y Proporciones}
\subtitle{$\chi^2$, $F$, Wald, Wilson y la Transformación de Fisher}
\author[J. Castillo Colmenares]{Juliho Castillo Colmenares}
\institute{Tecnológico de Monterrey}
\date{\vspace{-1.2cm}}

\begin{document}

% Slide 1: Portada
\begin{frame}[plain]
  \titlepage
\end{frame}

% Slide 2: Hoja de Ruta
\begin{frame}{Hoja de Ruta --- Capítulo 06: Estimación y su Relación con Ciencia de Datos}
  \footnotesize
  ¿Dónde nos encontramos en el desarrollo modular del capítulo? \pause
  \begin{itemize}\itemsep=0.08em
    \item \textbf{06.01} Estimación Puntual, Insesgadez, Eficiencia y Consistencia \pause
    \item \textbf{06.02} Método de Momentos (MoM) \pause
    \item \textbf{06.03} Estimación por Máxima Verosimilitud (MLE) y Score \pause
    \item \textbf{06.04} Intervalos de Confianza para Medias Poblacionales ($Z$ y $t$) \pause
    \item \textbf{06.05} Intervalos de Confianza para Varianzas y Proporciones \emph{(Hoy --- Cierre del Capítulo)}
  \end{itemize}
  \vspace{-0.05cm}
  \begin{block}{Objetivo de la Sesión}
    Cerrar el capítulo aplicando la misma lógica de estimación por intervalo --- estimador $\pm$ margen --- a la varianza poblacional (vía $\chi^2$), a la razón de dos varianzas (vía $F$), y a proporciones, incluyendo alternativas más robustas que el intervalo de Wald clásico.
  \end{block}
\end{frame}

% Slide 3: Motivación
\begin{frame}{Cerrando el Círculo: Todas las Distribuciones del Capítulo 05}
  \footnotesize
  \begin{motivacion}
    Hoy usamos, en un mismo mazo, las tres distribuciones muestrales que construimos en el Capítulo 05: $\chi^2$ para la varianza, $F$ para la razón de varianzas, y la Normal para proporciones. Es el broche perfecto para cerrar el Capítulo 06.
  \end{motivacion}

  \vspace{0.15cm}
  \pause
  \begin{itemize}\itemsep=0.1cm
    \item \textbf{Varianza:} $(n-1)S^2/\sigma^2 \sim \chi^2_{n-1}$ (Teorema de Fisher, Sección 05.03). \pause
    \item \textbf{Razón de varianzas:} cociente de chi-cuadradas escaladas $\sim F$ (Sección 05.05). \pause
    \item \textbf{Proporciones:} el estimador de Wald clásico tiene un competidor más robusto, el intervalo de Wilson.
  \end{itemize}
\end{frame}

% Slide 4: IC para Varianza y Razón de Varianzas
\begin{frame}{IC para $\sigma^2$ y para $\sigma_1^2/\sigma_2^2$}
  \footnotesize
  \begin{block}{Varianza Poblacional (Sección 05.03)}
    $$\frac{(n-1)S^2}{\chi^2_{\alpha/2,\,n-1}} \le \sigma^2 \le \frac{(n-1)S^2}{\chi^2_{1-\alpha/2,\,n-1}}$$
  \end{block}
  \pause

  \vspace{0.1cm}
  \begin{alertblock}{Razón de Dos Varianzas (Sección 05.05)}
    $$\frac{S_1^2/S_2^2}{F_{\alpha/2,\,n_1-1,\,n_2-1}} \le \frac{\sigma_1^2}{\sigma_2^2} \le \frac{S_1^2/S_2^2}{F_{1-\alpha/2,\,n_1-1,\,n_2-1}}$$
    Ambos intervalos son \textbf{asimétricos} alrededor del estimador puntual, a diferencia de los de $\mu$.
  \end{alertblock}
\end{frame}

% Slide 5: Proporciones: Wald y Wilson
\begin{frame}{Más Allá de Wald: el Intervalo de Wilson}
  \footnotesize
  \begin{block}{Intervalo de Wald (Clásico)}
    $$\hat p \pm Z_{\alpha/2}\sqrt{\hat p(1-\hat p)/n}$$
    Puede fallar gravemente para $n$ pequeña o $\hat p$ cerca de $0$/$1$.
  \end{block}
  \pause

  \vspace{0.1cm}
  \begin{alertblock}{Intervalo de Wilson (Score)}
    $$\frac{1}{1+\frac{Z_{\alpha/2}^2}{n}}\left[\hat p + \frac{Z_{\alpha/2}^2}{2n} \pm Z_{\alpha/2}\sqrt{\frac{\hat p(1-\hat p)}{n}+\frac{Z_{\alpha/2}^2}{4n^2}}\right]$$
    Se obtiene invirtiendo la prueba de score directamente; mantiene cobertura real mucho más cercana al nivel nominal.
  \end{alertblock}
\end{frame}

% Slide 6: Transformación de Fisher
\begin{frame}{Un Caso Especial: IC para la Correlación de Pearson}
  \footnotesize
  El coeficiente $r$ vive en $[-1,1]$: un IC simétrico ingenuo puede salirse de ese rango. \pause

  \vspace{0.1cm}
  \begin{block}{Transformación Estabilizadora de Varianza de Fisher}
    $$Z(r) = \text{artanh}(r) \sim N\!\left(\text{artanh}(\rho),\ \frac{1}{n-3}\right)$$
  \end{block}

  \vspace{0.1cm}
  \pause
  \begin{alertblock}{Construcción del IC}
    Se construye el IC simétrico para $Z(\rho)$ y se invierte con $\tanh(\cdot)$, produciendo un intervalo \emph{asimétrico} para $\rho$ que jamás excede $[-1,1]$.
  \end{alertblock}
\end{frame}

% Slide 7: Aplicaciones
\begin{frame}{Aplicaciones en Ciencia de Datos e Ingeniería}
  \footnotesize
  \begin{itemize}\itemsep=0.1cm
    \item \textbf{Control de calidad:} IC para la variabilidad de un proceso de manufactura. \pause
    \item \textbf{A/B testing:} comparación robusta de tasas de conversión, especialmente con muestras pequeñas. \pause
    \item \textbf{Finanzas cuantitativas:} IC para la volatilidad (varianza) de un activo. \pause
    \item \textbf{Análisis de correlación:} IC de Fisher para validar relaciones entre variables en modelos predictivos.
  \end{itemize}
\end{frame}

% Slide 8: Lab Python (Bloque 1)
\begin{frame}[fragile]{Laboratorio en Python: IC para Varianza y Wald vs. Wilson}
  \scriptsize
  IC $\chi^2$ para $\sigma^2$, y comparación Wald vs. Wilson para una proporción de muestra pequeña (\texttt{06.05\_confidence\_intervals\_variances.py}):
  {\lstset{basicstyle=\fontsize{4pt}{4.8pt}\selectfont\ttfamily,aboveskip=0.1em,belowskip=0.1em}
  \lstinputlisting[firstline=18, lastline=45]{../../code/06_estimacion_estadistica/06.05_confidence_intervals_variances.py}}
\end{frame}

% Slide 9: Lab Python (Bloque 2)
\begin{frame}[fragile]{Laboratorio en Python: Razón de Varianzas y Prueba A/B}
  \scriptsize
  IC $F$ para $\sigma_1^2/\sigma_2^2$, e IC para la diferencia de dos proporciones:
  {\lstset{basicstyle=\fontsize{4pt}{4.8pt}\selectfont\ttfamily,aboveskip=0.1em,belowskip=0.1em}
  \lstinputlisting[firstline=48, lastline=69]{../../code/06_estimacion_estadistica/06.05_confidence_intervals_variances.py}}
\end{frame}

% Slide 10: Lab Python (Bloque 3)
\begin{frame}[fragile]{Laboratorio en Python: Fisher y Cobertura Wald vs. Wilson}
  \scriptsize
  IC de Fisher para correlación, y verificación Monte Carlo de la cobertura real Wald vs. Wilson:
  {\lstset{basicstyle=\fontsize{3.6pt}{4.3pt}\selectfont\ttfamily,aboveskip=0.05em,belowskip=0.05em}
  \lstinputlisting[firstline=72, lastline=107]{../../code/06_estimacion_estadistica/06.05_confidence_intervals_variances.py}}
\end{frame}

% Slide 11: Salida Terminal
\begin{frame}[fragile]{Laboratorio en Python: Salida en Terminal}
  \scriptsize
  Salida estándar generada al ejecutar el script del laboratorio en Python:
  \vspace{0.02cm}
  \begin{block}{Salida Estándar en Consola}
    \fontsize{4pt}{5pt}\selectfont
    \begin{verbatim}
=== Block 1: Variance CI and Wald vs. Wilson ===
90% CI for sigma^2: [1.3387, 3.8753] | sigma: [1.1570, 1.9686]
Wald 95% CI: [0.0247, 0.3753] | Wilson 95% CI: [0.0807, 0.4160]

=== Block 2: Variance Ratio and A/B Test ===
98% CI for ratio: [0.7482, 10.9987] (contains 1? Yes)
95% CI for p1-p2: [-0.0785, -0.0015]

=== Block 3: Fisher Z-Transformation & Coverage ===
95% CI for rho: [0.6439, 0.9136]
Wald coverage: 0.8765 (nominal 0.95) | Wilson coverage: 0.9564
    \end{verbatim}
  \end{block}
\end{frame}

% Slide 12A: Ejercicio Nivel 1 (Enunciado)

% Slide 12B: Ejercicio Nivel 1 (Resolución)

% Slide 13A: Ejercicio Nivel 2 (Enunciado)

% Slide 13B: Ejercicio Nivel 2 (Resolución)

% Slide 14A: Ejercicio Nivel 3 (Enunciado)
\begin{frame}{Ejercicio en Clase --- Nivel 3: Analítico (1/2)}
  \footnotesize
  \begin{block}{Problema --- Razón de Varianzas (Enunciado)}
    Sensor 1: $n_1=16$, $S_1^2=0.36$°C². Sensor 2: $n_2=13$, $S_2^2=0.12$°C². Construya el IC del $98\%$ para $\sigma_1^2/\sigma_2^2$ y determine si hay evidencia de variabilidad distinta.
  \end{block}

  \vspace{0.1cm}
  \pause
  \begin{alertblock}{Estrategia}
    \scriptsize
    Use $F_{15,12,0.01}=3.96$ y $F_{15,12,0.99}=1/F_{12,15,0.01}\approx 0.2725$.
  \end{alertblock}
\end{frame}

% Slide 14B: Ejercicio Nivel 3 (Resolución)
\begin{frame}{Ejercicio en Clase --- Nivel 3: Analítico (2/2)}
  \scriptsize
  Calculamos el cociente empírico y el intervalo: \pause

  \vspace{0.05cm}
  \begin{block}{Cálculo}
    $$\frac{S_1^2}{S_2^2}=3.00, \qquad \text{IC}_{98\%}\left(\frac{\sigma_1^2}{\sigma_2^2}\right) = \left[\frac{3.00}{3.96},\,\frac{3.00}{0.2725}\right] = [0.758,\,11.01]$$
  \end{block}

  \vspace{0.05cm}
  \pause
  \begin{alertblock}{Interpretación}
    \scriptsize
    Como el intervalo contiene a $1$, no hay evidencia suficiente al $98\%$ de que un sensor sea más variable que el otro.
  \end{alertblock}
\end{frame}

% Slide 15A: Ejercicio Nivel 4 (Enunciado)
\begin{frame}{Ejercicio en Clase --- Nivel 4: Desafiante (1/2)}
  \footnotesize
  \begin{block}{Problema --- Transformación de Fisher (Enunciado)}
    Un estudio bioinformático con $n=28$ pacientes observa $r=0.82$. Calcule el IC del $95\%$ para $\rho$ usando la transformación de Fisher.
  \end{block}

  \vspace{0.15cm}
  \pause
  \begin{alertblock}{Estrategia}
    \scriptsize
    Calcule $Z(r)=\text{artanh}(r)$, construya el IC simétrico en la escala $Z$ con $\text{SE}=1/\sqrt{n-3}$, e invierta con $\tanh(\cdot)$.
  \end{alertblock}
\end{frame}

% Slide 15B: Ejercicio Nivel 4 (Resolución)
\begin{frame}{Ejercicio en Clase --- Nivel 4: Desafiante (2/2)}
  \scriptsize
  Transformamos, construimos el IC, e invertimos: \pause

  \vspace{0.05cm}
  \begin{block}{Cálculo}
    \scriptsize
    $$Z(0.82)=1.1568,\quad \text{SE}_Z=1/\sqrt{25}=0.2$$
    $$Z_L=0.7648,\ Z_U=1.5488 \implies \text{IC}_{95\%}(\rho) = [\tanh(0.7648),\,\tanh(1.5488)] = [0.644,\,0.914]$$
  \end{block}

  \vspace{0.05cm}
  \pause
  \begin{alertblock}{Interpretación}
    \scriptsize
    El intervalo es asimétrico alrededor de $r=0.82$ ($0.176$ hacia la izquierda vs. $0.094$ hacia la derecha) y nunca excede $[-1,1]$, a diferencia de lo que ocurriría con un IC ingenuo simétrico.
  \end{alertblock}
\end{frame}

% Slide 16: Cierre
\begin{frame}{Cierre de Sección: IC para Varianzas y Proporciones}
  \begin{reflexion}
    Con esta sección cerramos el Capítulo 06: desde definir qué hace ``bueno'' a un estimador (06.01), pasando por dos métodos de construcción (MoM, MLE), hasta usar ambos para construir intervalos de confianza exactos para medias, varianzas y proporciones. Cada intervalo comparte la misma lógica --- estimador $\pm$ margen --- calibrada por la distribución muestral correcta.
  \end{reflexion}

  \vspace{0.2cm}
  \pause
  \begin{block}{Lecciones Clave}
    \begin{itemize}\itemsep=0.08cm
      \item \textbf{Varianza:} IC vía $\chi^2$, asimétrico; \textbf{razón de varianzas:} IC vía $F$.
      \item \textbf{Proporciones:} Wilson supera a Wald en cobertura real, especialmente para $n$ pequeña.
      \item \textbf{Fisher:} transforma para estabilizar varianza y respetar los límites naturales del parámetro.
    \end{itemize}
  \end{block}
\end{frame}

% Slide 17: Perspectiva Modular
\begin{frame}{Perspectiva Modular --- Cierre del Capítulo 06}
  \scriptsize
  \begin{retoclase}
    Con la Sección 06.05 concluimos el Capítulo 06: Estimación y su Relación con Ciencia de Datos, dominando la estimación puntual, sus métodos de construcción, y la estimación por intervalo. El Capítulo 07 --- Docimasia (Pruebas de Hipótesis) --- formalizará la dualidad IC-prueba de hipótesis que hemos mencionado repetidamente: errores Tipo I/II, potencia estadística, y pruebas para una y dos muestras.
  \end{retoclase}

  \vspace{0.08cm}
  \pause
  \begin{alertblock}{Preparación Recomendada}
    \begin{itemize}\itemsep=0.06cm
      \item Repasa la dualidad IC--prueba de hipótesis mencionada en varios ejercicios de este capítulo.
      \item Reflexiona sobre qué significa rechazar o no rechazar $H_0$ cuando un valor hipotético cae dentro o fuera de un intervalo de confianza.
    \end{itemize}
  \end{alertblock}
\end{frame}

\end{document}
